% Part: applied-modal-logics
% Chapter: temporal-logic
% Section: introduction

\documentclass[../../../include/open-logic-section]{subfiles}

\begin{document}

\olfileid{aml}{tl}{int}

\olsection{Pambuka}

Logika temporal ngrembug pratelan ngenani bab kang bakal utawa wis
kedadeyan. Arthur Prior dianggep minangka pangrintis logika temporal,
kang diarani \emph{logika kala} dening dheweke. Pangrembugan logika
temporal ing kene umume manut cara modal asli Prior: operator temporal
dilebokake ing kerangka dhasar logika proposisional, kang lumrahe
nganggep pratelan tanpa kala.

Contone, ing logika proposisional, kita bisa ngrembug asu aran Beezie,
kang kadhang lungguh lan kadhang ora lungguh, kaya lumrahe asu.
Ing logika klasik, pratelan yen Beezie lagi lungguh lan uga ora lagi
lungguh bakal kontradiktif. Nanging, loro-lorone bisa bener ing
wektu kang beda; kanthi nambahake operator temporal ing basa,
kita bisa medharake bab iku kanthi gampang. Operator temporal uga
ngidini kita njlentrehake validitas inferensi saka ``Beezie bakal
nampa panganan cilik utawa bal'' dadi ``Beezie bakal nampa panganan
cilik utawa Beezie bakal nampa bal.''

Nanging, logika temporal nuwuhake akeh pitakon filsafat kang bisa
nyebabake kita milih kerangka logika temporal siji tinimbang liyane.
Contone, pratelan kontingen ngenani mangsa ngarep yaiku pratelan
ngenani bab kang durung kelakon, kang ora niscaya lan ora mokal.
Yen kita kandha ``Richard bakal lunga menyang toko panganan besuk,''
kita medharake pratelan ngenani bab kang durung kelakon lan kang
nilai bebeneré bisa didebat. Malah, bisa didebat apa pratelan iku
bisa \emph{diwenehi} nilai bebener wiwit wiwitan. Yen kita ngugemi
determinisme kang ketat, kita bisa nampa manawa ukara iku sejatiné
bener utawa salah sanajan kedadeyane durung tekan wektune---mung wae
kita durung ngerti nilai bebeneré. Kosok baline, kita bisa percaya
marang mangsa ngarep kang saestu mbukak, kanthi nilai bebener
pratelan kontingen ngenani mangsa ngarep durung ditemtokake.

Pranyata, akeh panemu ngenani struktur lan sipat wektu kasebut wis
kawangun ing pilihan model lan kerangka logika temporal kita.
Contone, kita bisa takon apa perlu mbangun model kang wektune
linier, nyabang, utawa malah siklis. Kita uga kudu mutusake apa
model temporal nduweni titik wiwitan lan pungkasan, lan apa wektu
bakal digambarake nganggo titik-titik diskret utawa minangka kontinum.

\end{document}
